\section{Introduction}
\label{sec:introduction}

Deciphering how cells transcriptionally respond to perturbations remains a central challenge in biology. Addressing it requires models that are predictive, scalable to the genome, and mechanistically interpretable~\cite{bunne2024build}. Mechanistic kinetic models do not scale, as measuring kinetic parameters experimentally is expensive, limiting them to simple systems~\cite{karr2012whole,hu2023comparative}. Deep neural network--based methods such as GEARS, scGPT, and STATE~\cite{lotfollahi2023predicting,roohani2024predicting,cui2024scgpt,levine2024cell2sentence,adduri2025predicting} scale to genome-wide prediction but require large training sets and offer little interpretability, and recent benchmarks show that they do not consistently surpass simple linear baselines~\cite{ahlmanneltze2025comparison}. Biologically informed neural networks (biNNs) aim to combine both via architectural priors, but existing biNNs mainly predict higher-level phenotypes~\cite{fortelny2020knowledge,lotfollahi2023biologically} rather than genome-wide expression responses to genetic perturbations.

In this paper, we introduce iPerturb, a biNN-based framework for predicting genome-wide expression changes induced by genetic perturbations, built on two ideas. As illustrated in Figure~\ref{fig:overview}, iPerturb takes a cell line and a genetic perturbation as input and works in three steps. First, cell-line context is used to query public regulatory resources and construct a \emph{context-specific gene regulatory network} (GRN) from five evidence tiers, keeping signed regulation and unsigned association evidence strictly apart. Second, the perturbation is propagated through the GRN by a \emph{learnable transfer function that follows the Hill function of gene regulation}~\cite{wittmann2009,eduati2017}, with the perturbed gene clamped at every update; because the Hill curve of every regulator is expressed in a Chebyshev basis, propagation becomes a matrix multiplication that runs efficiently on accelerators, a form related to basis-function expansions in learnable-activation networks~\cite{liu2024kan,li2024fastkan}. Third, every prediction is explained by the named neighbours and edges that support it, and the explanation is tested by ablation.

We applied iPerturb to CRISPRi perturbations in two human cell lines, K562 and RPE1, predicting the mean response of held-out target genes. Our results show that (i)~the context-specific GRN carries predictive signal for unseen perturbations that degree-preserving rewired and purely data-derived graphs do not, and that this signal survives the removal of response magnitude; (ii)~the Hill-form operator reproduces exact Hill propagation to within $10^{-5}$ while the matrix-multiplication form lowers memory six- to eight-fold on dense graphs and makes edge-specific learnable shapes feasible; and (iii)~on the curated signed core, which reaches only a small fraction of targets, Hill-form propagation matches but does not exceed linear propagation. Here, \emph{interpretable} is meant mechanistically: each predicted change decomposes into contributions from named neighbours and edges of the curated network, rather than from opaque latent factors; we do not claim that the explanation paths are validated molecular mechanisms.
